\clearpage
\appendix
\renewcommand{\thesection}{\Alph{section}}
\renewcommand{\sectionmark}[1]{\markright{Appendix \thesection\enspace #1}}

\section{Errors to fix whatever the paper becomes}\label{sec:app-errors}

These do not depend on which paper is written.

\subsection*{Statements that are wrong}
\begin{enumerate}
\item \textbf{Fact 5 and Section 8.5.} Report the event study relative to the last half-year before ChatGPT and drop the 7\%. The honest statement is that adopting the vocabulary did not raise output.
\item \textbf{Appendix B.8 and Fact 5.} Document the 2026 rise in papers per author. ``At most a few percent'' is false.
\item \textbf{Appendix B.8.} ``Decays with half-life $\ln2/\rho$'' is wrong under the paper's own law of motion: a shock of size $\alpha$ is absorbed in about $\ln(1+\alpha)/\rho$. And the appendix prices entry on the flow while Lemma~2 prices it on the stock.
\item \textbf{Abstract.} Remove ``identify the margins the model relies on'' and the range 1.3--7.2.
\item \textbf{Section 8.4.} A binding cap on writing would not be ``largely dissipated by entry''; by Corollary~2 it would pass through almost fully.
\item \textbf{Section 8.6.} The low volume elasticity does not come from the response of entry into science to its cost; it comes from the supply of new topics.
\item \textbf{The caption of the life-cycle figure.} With forward-looking entry the stopping density does not stay below the speed limit.
\end{enumerate}

\subsection*{Theory}
\begin{enumerate}[resume]
\item \textbf{Lemma 4 and Theorem 3.} One convention for how topics are born; the term for early cohorts; a numerical test that computes output per researcher from member counts.
\item \textbf{Lemma 3.} Take the founder's payoff along the path of her topic. Let founders leave topics that attract no follower. Call $\omega$ a share of time.
\item \textbf{Theorem 3.} Write the tool's effects $(h,f_0,m)$ explicitly instead of folding the change in the number of topics into $f$. State the sign of the pass-through on the falling side of $v$, where it can be negative.
\item \textbf{Corollary 2.} State it with fertility depletion, and call 0.99 an upper bound.
\item \textbf{Proposition 4(iii).} With no founders and no depletion the subsidised equilibrium is a tangency and violates the stability condition.
\item \textbf{Proposition 9.} Replace the two numerically verified steps by proofs. After a stall, with $\sigma=S/P$, $\dot\sigma=1-(\delta-\rho)\sigma>0$, and the value of joining falls along the path. For the threshold, prove the two-sided bound $\max\{g_0,g_{\hat P}\}\le g^*\le\max\{g_0,g^*_\infty\}$ (0.205 and 0.214 in the calibration). Rename the bound in part~(iii), which the text and the table confuse with the numerical threshold.
\item \textbf{Theorem 2(iv).} ``At least'' holds only at given reading shares; with shares responding the closure factor is 6.3, not 5.9.
\item \textbf{The life-cycle law of motion.} Add the catch-up delay.
\end{enumerate}

\subsection*{Calibration and tables}
\begin{enumerate}[resume]
\item Finish the parameter search on the fine grid, and report how the headline numbers move across near-optimal parameter vectors.
\item \textbf{The fit table.} Compare like with like: the model's regression peak (2.67) with the data's (1.0), and the share beyond each.
\item \textbf{The parameter table.} List the measurement-error constant (0.3). Correct the ``source'' of the erosion parameter and of the speed limit. Report members per topic.
\item \textbf{The incidence table.} Change one assumption per column; hold the share of explorers fixed when the founders' share varies.
\item \textbf{The reading share.} Define the model moment over all topic members, or target incumbents and newcomers separately.
\item Drop or reformulate the assumed premium target, which the calibrated curve cannot attain.
\item Report sensitivity to the assumed birth flow and to the decay rate of the literature.
\end{enumerate}

\subsection*{Text left over from the previous version}
\begin{enumerate}[resume]
\item Fact 3: the arXiv penalty and the scooping penalty of Hill and Stein ``validate each other'', against Section~2.3 (``they measure different events'').
\item Section 8.5 and two table notes: ``faster drafting'' attributed to a proposition about production speed, against Section~2.1 (``it is not drafting speed''); a ``drafting channel of Proposition~7'' that does not exist.
\item ``What we no longer claim'': channels move flagged overlap ``by 7--30\%''; the current table shows $+5$, $+2$, $+19$, $-11$, $-36$, $+35$.
\item Fact 1 cites prediction P7 for the radius pattern; it is P6. Prediction P11 cites P3 for newcomers' first papers; it is P4.
\item P6 is listed as out-of-sample although ``Fact 1 already shows this ordering''.
\item ``Five facts that discipline the model'': Fact~5 enters no parameter.
\item ``Clustering by area \ldots\ is part of the replication package'': it is not in the repository.
\item Appendix D: ``87\% at the 95th percentile'' is the figure for the 90th.
\item The README carries the previous title and says that reading, search and catch-up tools ``are competed away''.
\item Every sentence that narrates earlier versions of the paper.
\end{enumerate}

\subsection*{Notation}
\begin{enumerate}[resume]
\item At least thirteen symbols carry two or more meanings: $s$ (founding rate, subsidy, social output, time share); $\delta$ (decay rate, follow-on credit); $\tau$ (transfer share, stock tax, similarity threshold, tool index); $g$ (payoff function, diffusion rate); $a$ (production speed, topic age); $\omega$ (founders' share, age weights); $T$ (catch-up time, window length); $\alpha$ (shock size, weight of readership); $E$ and $\EE$; $A$ and $\mathcal A$; $D$ (two definitions); $\Psi$ (two proofs); $\nu$ and $\nu_0$; and $\lambda$ with five subscripts. Target: under twenty symbols in the main text.
\end{enumerate}

\section{Literature to add or to use more}\label{sec:app-lit}

Identified by the referee's reading, partly through a web search. The authors should verify each reference before citing it.

\begin{center}\small
\begin{tabularx}{\textwidth}{@{}P{2.45in}L@{}}
\toprule
Reference & Why \\
\midrule
Roback (1982); Harris and Todaro (1970); Gordon (1954); Moretti (2011); Kline and Moretti (2014) & The incidence theorem is the free-entry incidence result of spatial and labour economics, and the insiders' premium is a queue rent. State this and claim the application. \\
Hill, Yin, Stein, Wang, Wang and Jones (\emph{Nature}, 2025), ``The pivot penalty in research'' & Direct evidence on the cost of entering a topic. \\
Gans, ``A Quest for AI Knowledge'' (NBER working paper) & Artificial intelligence and the direction of research. \\
Agrawal, McHale and Oettl, ``AI in Science'' (NBER working paper) & The closest economics of AI as a research tool. \\
Hoelzemann, Manso, Nagaraj and Tranchero, ``The streetlight effect in data-driven exploration'' (NBER working paper) & Close to ``reading tools concentrate science''. \\
Baccara, Levy and Razin, ``Research Waves'' & Life cycles of topics with forward-looking entry. \\
Bergstrom, Foster and Song (2016), ``Why scientists chase big problems'' & Crowding of problems under priority credit. \\
Teodoridis, Bikard and Vakili (2019) & Fast- and slow-paced domains of knowledge. \\
Models of information congestion (Van Zandt 2004; Anderson and de Palma 2009; Falkinger 2008) & ``The crowders are the readers''. \\
\midrule
\multicolumn{2}{@{}l}{\emph{Already in the bibliography, but cited only in a list or not at all in the text}} \\
\midrule
Hopenhayn and Squintani (2021); Bryan and Lemus (2017) & Already have free-entry overcrowding of hot research lines. \\
Klepper (1996); Jovanovic and MacDonald (1994) & The shakeout analogue of take-off and overshooting. \\
Azoulay et al.\ (2019) & The rival account of incumbents' rents. \\
Kusumegi et al.\ (2025); Hao et al.\ (2026) & The 2026 fact and ``speed thickens'' must be positioned against both. \\
\bottomrule
\end{tabularx}
\end{center}

\section{Where things are}\label{sec:app-files}

\begin{center}\small
\begin{tabularx}{\textwidth}{@{}P{2.3in}L@{}}
\toprule
Path in the repository & Content \\
\midrule
\texttt{paper/main.pdf}, \texttt{paper/sections/} & The draft discussed here (commit \texttt{866dcfd}). \\
\texttt{versions/v2\_2026-10-23\_\ldots/} & The previous version, complete. \\
\texttt{code/ftmodel/}, \texttt{code/tests/} & The model and its 111 numerical tests. \\
\texttt{code/results/*.json} & Calibrated parameters and every reported number. \\
\texttt{data/}, \texttt{empirics/} & The arXiv snapshot, the high-energy physics pairs, and the scripts and results of the six designs. \\
\texttt{docs/paper\_diagnosis.md} & The diagnosis and recommendation in short form. \\
\texttt{docs/theory\_revision.md} & What changed between versions 2 and 3, comment by comment. \\
\texttt{handout/} & This handout. \texttt{python3 handout/make\_figures.py} rebuilds its charts from the data and the model code; \texttt{latexmk -pdf handout.tex} builds the PDF. \\
\bottomrule
\end{tabularx}
\end{center}
